\documentclass{article}
\usepackage[UTF8]{ctex}
\usepackage{amsmath,amssymb,amsfonts}
\usepackage{enumitem}
\usepackage{geometry}
\geometry{a4paper, margin=1in}
% === 错题本排版优化（不参与题目标识） ===
\setlength{\parindent}{0pt}
\setlength{\parskip}{0.35em plus 0.10em minus 0.08em}
\setlength{\emergencystretch}{3em}
\setlist{itemsep=0.25em, topsep=0.40em, parsep=0pt}
\allowdisplaybreaks[2]
% === 错题本排版优化结束 ===
\begin{document}
\section*{中南大学2024年数学分析真题}
\section*{144．中南大学2024年数学分析试卷}

\subsection{题目1}
1 ．求极限 \(\lim _{n \rightarrow \infty}\left(1-\frac{3}{4 n^{2}}\right)^{n^{2}+4 n+3}\)

\subsection{题目2}
2．已知 \(f(x)=\left|\begin{array}{ccc}e^{x} & x^{2} & 2 \\ 1 & 3 x & 4 x^{2} \\ 0 & 5 & e^{2 x}\end{array}\right|\) ，求 \(f^{\prime}(x)\) ．

\subsection{题目3}
3．求定积分 \(\int_{0}^{\frac{\pi}{2}} \frac{\cos \theta}{\cos \theta+\sin \theta} \mathrm{d} \theta\) ．

\subsection{题目4}
4．求级数 \(\sum_{n=1}^{\infty} \frac{1}{n!}\) 的和。

二．（10 分）设 \(\left\{f_{n}(x)\right\}\) 在 \([a, b]\) 上满足：对任意的 \(x \in[a, b], f_{n}^{\prime}(x)\) 存在．且\\
（1）对任意的 \(x_{0} \in[a, b],\left\{f_{n}\left(x_{0}\right)\right\}\) 收敛。\\
（2）\(\left\{f_{n}^{\prime}(x)\right\}\) 在 \([a, b]\) 上一致收敛．\\
证明：\(\left\{f_{n}(x)\right\}\) 在 \([a, b]\) 上一致收敛．

三．（20 分）证明： \(\int_{0}^{1} \frac{\ln (1-t)}{t} \mathrm{~d} t=-\sum_{n=1}^{\infty} \frac{1}{n^{2}}\) ．

四．（15 分）设 \(f(x)\) 在 \([0, \pi]\) 上二阶连续可导，\(f(\pi)=2\) ，且

\[
\int_{0}^{\pi}\left[f(x)+f^{\prime \prime}(x)\right] \sin x \mathrm{~d} x=5
\]

求 \(f(0)\) ．

五．（15 分）计算以下曲线积分

\[
\int_{C} x^{2} y \mathrm{~d} x+(y-3) \mathrm{d} y
\]

其中 \(C\) 为长方形区域的边界，沿顺时针方向，长方形四个顶点分别为（0，0），（0，1），（2，0），（2，1）．

六．（15 分）定义以下二元函数

\[
f: \mathbb{R}^{2} \rightarrow \mathbb{R}, f(x, y)= \begin{cases}\frac{x^{3}+x^{2} y+x y^{2}+2 y^{4}}{x^{2}+y^{2}}, & x^{2}+y^{2} \neq 0 \\ 0, & x^{2}+y^{2}=0\end{cases}
\]

证明：

\begin{itemize}
  \item[（1）] \(f\) 在 \((0,0)\) 点连续．
  \item[（2）] \(f\) 在 \((0,0)\) 点处存在偏导数．
  \item[（3）] \(f\) 在 \((0,0)\) 点处是否可微？请详细说明．
\end{itemize}

七．（15 分）设 \(f:(0,1) \rightarrow \mathbb{R}\) 一阶连续可导，且满足对任意的 \(x, y \in(0,1)\) ，有 \(|f(x)-f(y)| \leq|x-y|^{2}\) ，证明：\(f\) 恒为常数．

八．（20 分）已知 \(f(x)\) 是 \((0,+\infty)\) 上的单调连续函数，且 \(\lim _{x \rightarrow+\infty} f(x)=0\) ．证明：

\[
\lim _{n \rightarrow \infty} \int_{0}^{+\infty} f(x) \sin (n x) \mathrm{d} x=0
\]
\end{document}
